%==========================================================================
% Exemplo de artigo cientifico UNESC
%
% v0.0 versão em 20/09/2023 por Marlon Oliveira
% v1.0 versão em 25/06/2025 por Marlon Oliveira


%==========================================================================

%==========================================================================
% Pacotes utilizados pela UNESC - não delete as linhas abaixo
%==========================================================================
\documentclass{UNESC}
\usepackage[T1]{fontenc}
\usepackage[utf8]{inputenc}%\usepackage[latin1]{inputenc}
\renewcommand{\rmdefault}{phv} % Arial
\renewcommand{\sfdefault}{phv} % Arial
\usepackage[brazilian]{babel}
\usepackage{graphicx}
%\usepackage{subfigure}
\usepackage{hyperref}
\usepackage{quoting}
\usepackage[labelsep=none]{caption}
\DeclareCaptionLabelFormat{figuracomtraco}{#1~#2 -}
\captionsetup[figure]{labelformat=figuracomtraco, labelsep=space}
\captionsetup[table]{labelformat=simple, labelsep=space, textformat=simple}
\DeclareCaptionLabelFormat{tabelacomtraco}{#1~#2 -}
\captionsetup[table]{labelformat=tabelacomtraco}
\captionsetup{labelsep=none}
%\renewcommand{\thefigure}{\arabic{figure} - }
\usepackage{enumitem}
\usepackage{float}
\usepackage{comment}
\usepackage[hang,flushmargin]{footmisc}
\usepackage{array}

\hypersetup{
  setpagesize  = false,
  colorlinks   = true,    % Colours links instead of ugly boxes
  urlcolor     = blue,    % Colour for external hyperlinks
  linkcolor    = black,   % Colour of internal links
  citecolor    = black    % Colour of citations
}
\urlstyle{same}
\usepackage{setspace}

\usepackage[alf,recuo=0.0cm,abnt-emphasize=bf]{abntex2cite} 

%\captionsetup[figure]{justification=raggedright}


%==========================================================================
% Dados do artigo: título, autores e filiação
%==========================================================================
\title{Monitoramento Inteligente de Sujidade:
Detecção de perdas por sujidade através do desvio entre curvas de potência real e teórica.}
\author
{
   Matheus Luís Mantovani Passoni\footnote{Autor, Curso de Ciência da Computação, Universidade do Extremo Sul Catarinense (Unesc), Criciúma - Santa Catarina - Brasil. matheus.mantovani28@gmail.com},
   %Nome Completo do corientador(a)\footnote{Corientador(a), Curso de Ciência da Computação, Universidade do Extremo Sul Catarinense (Unesc), Criciúma - Santa Catarina - Brasil. email@unesc.net}
   Allan Farias Fávaro\footnote{Orientador(a), Curso de Ciência da Computação, Universidade do Extremo Sul Catarinense (Unesc), Criciúma - Santa Catarina - Brasil. email@unesc.net}
}
%==========================================================================
% Início do texto do artigo
%==========================================================================
\begin{document}

\onehalfspacing

\maketitle


\section{INTRODUÇÃO}

A crescente demanda por fontes de energia renovável tem impulsionado o uso de sistemas fotovoltaicos como alternativa sustentável para geração de eletricidade. No Brasil, esse crescimento é expressivo: em 2023, a fonte solar correspondeu a 16,1\% da matriz elétrica do país. Desse total, a maior parte é proveniente de sistemas de geração distribuída e microgeração, com forte presença em residências e comércios \cite{Tramontin_Rampinelli_2024}. No entanto, diversos fatores ambientais influenciam diretamente o desempenho desses sistemas, destacando-se o acúmulo de poeira sobre os módulos solares, fenômeno conhecido como \textit{soiling}.

O acúmulo de partículas na superfície dos painéis reduz a incidência de radiação solar, impactando negativamente a geração de energia \cite{Araujo_Marques}. Estudos apontam que, após fatores como irradiância e temperatura, o \textit{soiling} é um dos principais responsáveis pela perda de eficiência em sistemas fotovoltaicos \cite{Tramontin_Rampinelli_2024}, podendo causar reduções significativas no desempenho dependendo das condições ambientais e do tipo de material particulado \cite{john2016, jaszczur2020}. Avaliações práticas em instalações reais demonstram que, em períodos sem chuva de apenas 17 dias, a perda de eficiência de uma usina pode chegar a 16,52\%, resultando em uma redução de 11,71\% na eletricidade gerada \cite{Araujo_Marques}. Além disso, em ambientes urbanos ou industriais, compostos da sujeira podem amplificar essas perdas, que se tornam ainda mais severas em condições de baixa irradiância solar \cite{Tramontin_Rampinelli_2024}.

Somado a isso, a variabilidade de fatores como localização geográfica, condições climáticas e composição da poeira torna a estimativa dessas perdas uma tarefa complexa \cite{Araujo_Marques}. Métodos tradicionais de análise ainda apresentam limitações, sendo frequentemente baseados em abordagens empíricas ou pouco precisas \cite{abuzaid2022}. No contexto da microgeração distribuída, essa limitação é ainda mais crítica, pois a manutenção e a limpeza dos módulos raramente seguem um rigor técnico, levando à subutilização crônica da capacidade instalada.

Diante desse cenário, surge a necessidade de soluções computacionais capazes de monitorar e analisar o desempenho dos sistemas fotovoltaicos de forma inteligente. A comparação entre a potência real gerada e a potência teórica esperada se apresenta como uma abordagem promissora para identificação de desvios associados à perda de eficiência, uma vez que anomalias na curva de operação elétrica são assinaturas claras de problemas de funcionamento \cite{Araujo_Marques}.

Nesse contexto, este trabalho propõe uma solução baseada no monitoramento inteligente da eficiência de sistemas fotovoltaicos, focada na detecção de perdas causadas por \textit{soiling} por meio da análise do desvio entre as curvas de potência real e teórica. Para implementar e validar essa abordagem, será utilizada uma ferramenta computacional capaz de analisar os dados de operação do sistema, comparando a energia efetivamente gerada com os modelos de capacidade esperada. A avaliação contínua desse desvio de desempenho tem se mostrado uma metodologia altamente eficaz na literatura recente para isolar e quantificar o impacto da sujidade com precisão \cite{Tramontin_Rampinelli_2024}. Ao analisar o desempenho teórico em contrapartida ao real, a solução proposta fornecerá diagnósticos automatizados, substituindo estimativas imprecisas por um modelo embasado que auxiliará na tomada de decisão para a limpeza otimizada dos painéis, mitigando as perdas operacionais na geração distribuída.

\subsection{Objetivo Geral}

Propor um método para monitoramento da eficiência de sistemas fotovoltaicos, com foco na identificação de perdas de desempenho por sujidade a partir da análise do desvio entre potência real e teórica.

\subsection{Objetivos Específicos}

\begin{itemize}

\item Coletar e pré-processar os dados operacionais de sistemas fotovoltaicos, estruturando a base necessária para a comparação analítica;

\item Analisar dados reais de geração de energia para avaliar o comportamento operacional e detectar perdas de desempenho;

\item Identificar padrões de desvio entre as curvas de potência real e teórica que sejam associados ao acúmulo de poeira;

\item Desenvolver um sistema computacional para automatizar o monitoramento e a detecção de anomalias na eficiência;

\item Validar a solução proposta por meio de testes computacionais, confrontando os diagnósticos do sistema com o histórico real de operação.

\end{itemize}

\subsection{Justificativa}

A eficiência dos sistemas fotovoltaicos é um fator crítico para a viabilidade econômica e operacional da geração de energia solar. Entre os diversos fatores que impactam essa eficiência, o acúmulo de poeira se destaca como um dos principais responsáveis pela degradação do desempenho dos painéis, resultando em perdas expressivas de energia ao longo do tempo \cite{abuzaid2022, Araujo_Marques}. A quantificação exata dessa perda de potência é fundamental, pois o soiling não apenas reduz a geração diária, mas altera os parâmetros elétricos de operação do módulo, diminuindo a corrente gerada e comprometendo a rentabilidade projetada para o sistema \cite{Araujo_Marques}.

A literatura aponta que a análise do impacto do soiling é complexa devido à influência de múltiplos fatores ambientais, como temperatura, umidade, vento e as propriedades físicas e químicas do material particulado, além da interação entre esses elementos \cite{abuzaid2022, jaszczur2020, john2016}. Mais do que o simples bloqueio da radiação, é imperativo monitorar esse acúmulo porque a deposição desigual de sujeira pode provocar o superaquecimento localizado das células fotovoltaicas, fenômeno que acelera a degradação física dos componentes e eleva o risco de falhas irreversíveis no sistema \cite{jaszczur2020}. Isso evidencia a necessidade crítica de abordagens mais robustas e automatizadas para o monitoramento contínuo da eficiência.

Do ponto de vista computacional, há uma lacuna na aplicação de soluções inteligentes capazes de identificar, de forma automatizada, as perdas de desempenho estritamente associadas ao acúmulo de sujeira, visto que muitas manutenções ainda dependem de cronogramas fixos ou inspeções visuais \cite{Tramontin_Rampinelli_2024}. O desenvolvimento de um modelo analítico com essa finalidade contribui tanto para o avanço acadêmico, ao propor novas metodologias de detecção de anomalias, quanto para aplicações práticas no setor energético.

Além disso, a utilização de abordagens computacionais para a análise de desempenho auxiliará diretamente na otimização de estratégias de manutenção. Ao verificar a perda de potência por sujidade com precisão, torna-se possível identificar o ponto de equilíbrio econômico entre o custo da limpeza física dos módulos e o prejuízo financeiro causado pela energia não gerada \cite{Araujo_Marques}. Dessa forma, a solução proposta fornece um diagnóstico preciso do sistema, reduzindo custos operacionais desnecessários e maximizando a eficiência energética das instalações.


%-----------------------------------------------
% Comentar para o Projeto em diante.
%-----------------------------------------------
\section{\small ASSINATURAS E TERMO DE ORIENTAÇÃO}
%-----------------------------------------------

\noindent\begin{tabular}{ | m{25em}  | m{2.5cm} | } 
  \hline
  ALUNO: &DATA: \\
  & \\ 
  & \\ 
  \hline
  PROFESSOR ORIENTADOR:&DATA:  \\
  & \\ 
  & \\ 
  \hline
  PROFESSOR CO-ORIENTADOR: &DATA:\\
  & \\ 
  & \\ 
  \hline
\end{tabular}


\bibliographystyle{abntex2-unesc}
\renewcommand{\refname}{REFERÊNCIAS}
\bibliography{referencias} % arquivo referencias.bib que contém as referências bibliográficas utilizadas no artigo


\end{document} 

